\section{What can a finite torus tell us about Yang--Mills theory?}\label{ym:sec:finite-guide}

Parallel transport around a small square measures curvature. Parallel transport
around a spatial circle can create a winding flux tube. Propagation through a
long Euclidean time interval measures excitation energies. On a periodic
lattice all three operations use the same link variables, but they answer
different questions. The distinction matters if one wants to infer a property
of an infinite-volume Hamiltonian from a calculation on a finite box.

There are two useful comparisons of partition functions. A \emph{twist}
changes the flux sector while keeping the box fixed. A \emph{doubling} changes
the box while keeping the local action, the bare coupling and the untwisted
boundary conditions fixed. The first measures the cost of imposing global
flux. The second, if it is sufficiently close to an extensive-volume law,
can certify suppression of \emph{all} excited transfer states. There is no
known implication here from a small twist free energy to a sufficiently small
doubling defect. A finite model in which both can be calculated exactly makes
the difference visible.

\subsection{From a connection to a probability law on links}
Put a cubical cell structure on a four-torus with $n_\mu$ sites in direction
$\mu$ and spacing $a$; the physical lengths are $a n_\mu$. A positive link
$(x,\mu)$ carries a parallel transporter $U_\mu(x)\in G$. Reversing a link
uses its inverse, not a second independent random variable. An elementary
square is a \emph{plaquette}. Its holonomy is the ordered product
\begin{equation}\label{ym:guide-holonomy}
 U_p=U_\mu(x)U_\nu(x+e_\mu)
 U_\mu(x+e_\nu)^{-1}U_\nu(x)^{-1}.
\end{equation}
Changing the frame at every vertex conjugates $U_p$ at the base point, so its
trace does not depend on the frame. For a smooth connection,
Eq.~\eqref{ym:plaquette-expansion} gives $U_p=\exp(a^2F_{\mu\nu}+O(a^3))$.
Thus $1-\operatorname{Re}\operatorname{tr} U_p/m$, for $G\subset U(m)$,
measures squared curvature to leading order, by Eq.~\eqref{ym:plaquette-energy}.
This is the \emph{classical} reason for the local action; it does not say
that typical configurations in the quantum integral are smooth connections.

Write $E$ for the sum of these dimensionless plaquette energies, each based
square in each of the six planes counted once. Integrate $e^{-\beta E}$ over
independent normalized Haar measures on the positive links, as in
Eq.~\eqref{ym:wilson-partition}. Dividing by this integral $Z$ gives the
Wilson probability law. At $\beta=0$ all assignments have their original Haar
weight and $Z=1$; at large $\beta$ the law favours small plaquette energy.
Small $\beta$ is \emph{strong} bare coupling, not a fine-mesh approximation
to a smooth field. A local observable accessible in a simulation is the mean
normalized plaquette trace. One cannot read a mass from its value on one box.

\subsection{Changing the flux: a worked example in cellular cohomology}
Take $G=\mathbb Z_2=\{+1,-1\}$. Write $U_e=(-1)^{u_e}$, so $u$ is a
cellular $1$-cochain over $\mathbb F_2$. A vertex gauge change adds the
coboundary of a $0$-cochain. Equation~\eqref{ym:guide-holonomy} now says
$U_p=(-1)^{(\delta u)_p}$: the plaquette field is an \emph{exact}
$2$-cochain. In particular $\delta^2u=0$; around each cube the product
of the six face signs is $+1$. The dual plaquettes of negative sign form
a closed sheet modulo two, not necessarily a smooth or connected surface.
Links are indexed by site and direction even on a side of length one;
repeated occurrences in a boundary cancel in pairs. Replacing the cell
structure by a simple graph would change the theory.

To impose flux, prescribe a closed $2$-cochain $t$ and replace $\delta u$
by $t+\delta u$ in the weight. Moving the chosen sheet by $\delta b$
changes no partition function: $t+\delta b+\delta u=t+\delta(u+b)$, and
$u\mapsto u+b$ permutes the integration variables. What cannot be changed
this way is the pairing of $t$ with a closed coordinate two-torus $S$:
\begin{equation}\label{ym:guide-flux-pairing}
 \langle t+\delta b,S\rangle=\langle t,S\rangle+\langle b,\partial S\rangle
 =\langle t,S\rangle\quad\text{in }\mathbb F_2.
\end{equation}
For example, mark the $12$-plaquettes based at $x^1=x^2=0$ for every
$x^3,x^4$. Each $12$-torus at fixed $(x^3,x^4)$ meets the mark once. The
marked plaquettes are dual to a sheet extending in directions $3,4$; its
location is movable but its flux is not. The six independent plane pairings
give $H^2(T^4;\mathbb F_2)\cong\mathbb F_2^6$, hence 64 twist sectors.
This is a background flux, not the result of gauge fixing a sampled field.

One can sum by plaquette patterns instead of link assignments. Let
$C=\operatorname{im}(\delta:C^1\to C^2)$ have rank $d$, let $w(c)$ count the
nonzero faces of $c$, and let $A_w^{(t)}$ count patterns of weight $w$ in
$t+C$. With $N_1$ links there are $2^{N_1-d}$ preimages of each pattern.
Each negative plaquette costs $1-(-1)=2$. The normalized integral is therefore
\begin{align}
 Z_\beta[t]&=2^{-N_1}\sum_u e^{-2\beta w(t+\delta u)}\notag\\
 &=2^{-d}\sum_w A_w^{(t)}e^{-2\beta w}.\label{ym:guide-enumerator}
\end{align}
This is an exact finite sum at any $\beta$, not a sample of rare flux events.
There is a check across sectors: the binary vectors orthogonal to $C$ are
$2$-cycles since $\langle z,\delta u\rangle=\langle\partial z,u\rangle$.
On the dual four-torus they are closed $2$-cochains, the disjoint union
of all 64 cohomology classes. The check takes a simple polynomial form. For a set $B$ of patterns on
$P$ faces put $W_B(X,Y)=\sum_{z\in B}X^{P-w(z)}Y^{w(z)}$. Orthogonality
gives $\mathbf 1_{z\in C^\perp}=2^{-d}\sum_{c\in C}(-1)^{c\cdot z}$.
In the sum over $z$, a face with $c_p=0$ then gives $X+Y$, and one with
$c_p=1$ gives $X-Y$. Identifying the dual torus with the original yields
\begin{equation}\label{ym:guide-duality}
 \sum_{[t]\in H^2(T^4;\mathbb F_2)} W_{t+C}(X,Y)
 = W_{C^\perp}(X,Y)=2^{-d}W_C(X+Y,X-Y).
\end{equation}
On the $2^4$ torus there are 96 faces and $d=45$; the two sides agree in
all 97 integer coefficients after clearing $2^d$. This tests the
\emph{sum} over sectors, not each sector against compensating errors.

For one magnetic twist the dimensionless free-energy cost is
$-\log(Z_\beta[t]/Z_\beta[0])$ (Eq.~\eqref{ym:twist-free-energy}). Choosing
a time direction, an electric flux sector is selected instead by Fourier
projection over the three temporal twist bits $k$ at zero magnetic twist:
\begin{equation}\label{ym:guide-electric}
 \widehat Z_\beta(e)=\tfrac18\sum_{k\in\mathbb F_2^3}(-1)^{e\cdot k}Z_\beta[k],
 \qquad F_e=-\log\frac{\widehat Z_\beta(e)}{\widehat Z_\beta(0)} .
\end{equation}
For example, on $2^4$ at $\beta=0.2$, Eq.~\eqref{ym:guide-enumerator} gives
a magnetic cost of about $0.013$ and Eq.~\eqref{ym:guide-electric} an electric
cost of about $5.06$ for one unit; at $\beta=0.6$ the costs are about $3.21$
and $0.086$, respectively. The relative cost reverses on the \emph{same}
small torus. Confinement is a question about how such costs, or large Wilson
loops, scale with size; a finite-box crossing does not settle it. A centre
twist is available also in $SU(2)$, but an $SU(2)$ link is not a binary
cochain and Eq.~\eqref{ym:guide-enumerator} does not apply to that group.

\subsection{Changing the box: why a seed says something about a spectrum}
Keep the transverse sizes fixed and call one direction time. Collecting the
spatial links into a slice and integrating the temporal links gives a
one-step operator $\mathsf T$. Closing a history after $t$ slices gives
$Z(t)=\operatorname{Tr}\mathsf T^t$ (Eq.~\eqref{ym:trace-formula}). It is
operator positivity, not just a pointwise positive Boltzmann factor, that
lets us interpret this trace as a sum of nonnegative powers. In the sign
model the identity $e^{\beta s}=\cosh\beta+s\sinh\beta$ for $s=\pm1$
expands the temporal weight in products of link signs. Averaging temporal
links kills a product odd at any vertex. Each surviving term is a
nonnegative coefficient times the same slice function on the two ends
(an outer product), so $\mathsf T$ is positive semidefinite. For compact
matrix groups the corresponding argument is Proposition~\ref{ym:sc-prop:positivity}.

Here is what a time doubling measures. If $\lambda_j$ are the eigenvalues
with multiplicity, form the normalized weights $p_j=\lambda_j^t/Z(t)$.
Then $\sum_j p_j=1$ and
\begin{equation}\label{ym:guide-purity}
 \frac{Z(2t)}{Z(t)^2}=\sum_j p_j^2\le p_1\sum_j p_j=p_1.
\end{equation}
Purity near one forces one state to dominate: if $\mathsf P$ denotes the
left side, $\sum_{j>1}(\lambda_j/\lambda_1)^t=(1-p_1)/p_1\le
\mathsf P^{-1}-1$. No choice of interpolating field was made. Unlike a
single correlation measurement, the bound includes states to which a
chosen observable might not couple at all.

Now double \emph{all four} sides. The defect
$\Delta_\beta(n)=16\log Z_\beta(n^4)-\log Z_\beta((2n)^4)$ compares
boxes at the same $\beta$; here $n^4$ denotes the shape, not a time extent.
A contribution to $\log Z$ proportional to the number of sites cancels.
For one side-doubling the cost is $-\log\mathsf P\ge0$ by
Eq.~\eqref{ym:guide-purity}. Doubling successive axes expresses the defect
as the sum of four such costs with weights $8,4,2,1$
(Eq.~\eqref{ym:defect-telescope}). Thus a small defect cannot conceal
a large cost in one direction by cancellation against the others.

A \emph{seed} is an inequality at one starting size, $\Delta_\beta(n)\le
1/7480$, not an initial field configuration. What makes that particular
test useful is Theorem~\ref{ym:thm:recurrence}. Put
$h_k=3740\,\Delta_\beta(2^kn)$. In the seed range it gives
\begin{equation}\label{ym:guide-iterate}
 h_0\le\tfrac12,\qquad h_{k+1}\le h_k^2,
 \qquad h_k\le(\tfrac12)^{2^k}.
\end{equation}
The speed of this decay matters: the time interval in the purity test
also doubles. Take time as the \emph{first} doubled side of a cube of side
$N=2^kn$. Its transverse side is still $N$. Equations~\eqref{ym:guide-purity}
and \eqref{ym:defect-telescope} give, for $r_N=\lambda_2/\lambda_1$,
\begin{equation}\label{ym:guide-first-axis}
 r_N^N\le\mathsf P^{-1}-1\le e^{\Delta_\beta(N)/8}-1.
\end{equation}
Since $e^x-1\le2x$ for $0\le x\le1$, Eq.~\eqref{ym:guide-iterate}
bounds the right side of Eq.~\eqref{ym:guide-first-axis} by
$2^{-2^k}/(4\cdot3740)$. Taking the $2^k$-th root gives $r_N^n\le1/2$.
The exponent in this last bound is the \emph{starting} side $n$ at every
volume. Choosing time last instead gives the $2^{k+1}n$ transverse side
in Eq.~\eqref{ym:conditional-gap}.

\paragraph{A seed one can actually check.}
On the $1^4$ sign torus every plaquette is a commutator of signs, hence $+1$
and $Z_\beta(1^4)=1$. On $2^4$ the four links of each plaquette are distinct;
at $\beta=0$ their signs are independent and unbiased. Consequently each of
96 faces is negative with probability $1/2$ (the faces need not be independent).
In Eq.~\eqref{ym:guide-enumerator}, $\langle w\rangle_0=48$. It follows that
\begin{equation}\label{ym:guide-first-seed}
 \Delta_0(1)=0,\qquad
 \left.\frac{\partial\Delta_\beta(1)}{\partial\beta}\right|_{\beta=0}
 =-\left.\frac{\partial\log Z_\beta(2^4)}{\partial\beta}\right|_{\beta=0}
 =2\langle w\rangle_0=96.
\end{equation}
Evaluating the same finite sum gives the crossing $\Delta_\beta(1)=1/7480$
at $\beta\simeq1.3926034652\times10^{-6}$. The much cheaper sufficient
bound $\Delta_\beta(n)\le192\beta n^4$ of Eq.~\eqref{ym:strong-seed}
certifies $\beta\le1/(7480\cdot192)$ at $n=1$. The bound has exactly twice
the slope in Eq.~\eqref{ym:guide-first-seed}; the positive crossing
couplings are only \emph{approximately} in that ratio.

For the four-dimensional sign model the whole finite statement has an
independently checked formal proof: for any integer $n\ge1$, a single
$0\le\beta\le 1/(7480\cdot192\,n^4)$, every time direction and every
$k\ge0$, $r_{2^kn}^{\,n}\le1/2$. The trace representation, positivity and
seed are derived from the lattice rather than supplied as hypotheses.
If $r_N>0$ this gives $-\log r_N\ge(\log2)/n$ in lattice units; if
$r_N=0$ the ratio statement already applies without taking a logarithm.
One can see which eigenvalue this means without diagonalizing a large
matrix. On a spatial slice of side one there are three self-links. Their
spatial plaquettes are trivial; the two occurrences of the temporal sign
in each time-like plaquette cancel. If a spatial sign is unchanged on
the next slice its weight is 1; if it flips its weight is
$q=e^{-2\beta}$. With normalized counting measure the one-link matrix is
\begin{equation}\label{ym:guide-one-link}
 B=\tfrac12\begin{pmatrix}1&q\\q&1\end{pmatrix},\qquad
 B\binom{1}{\pm1}=\frac{1\pm q}{2}\binom{1}{\pm1}.
\end{equation}
The full matrix is $B^{\otimes3}$: replacing one of the three even
eigenvectors by the odd one gives $r_1=(1-q)/(1+q)=\tanh\beta$.
At $\beta=0$ only the constant state propagates. At positive $\beta$ the
second mode is real and nonzero. This is a control in a known
strong-coupling regime, not a continuum result for a compact simple Lie group.

\subsection{What is measured, and what must survive a limit?}
For $SU(2)$ it is easier to sample a plaquette than to determine an absolute
partition function. Equation~\eqref{ym:num-integration} integrates the
\emph{difference} of mean plaquettes on two boxes to obtain the defect:
the normalization is fixed by $Z_0=1$. If a scan starts instead at
$\beta_0>0$, it measures $\Delta_\beta-\Delta_{\beta_0}$; an increment
compatible with zero is not an upper bound on $\Delta_\beta$. Moreover,
a seed at $n$ requires $n\,\mathsf m(n)\gtrsim 11$ by
Eq.~\eqref{ym:num-necessary}. A fitted mass is not a certified upper bound
on the lowest excitation, so this condition cannot be checked by simply
inserting a fitted number.

A contractible spatial loop summed over its position probes a scalar
glueball channel when projected onto the scalar symmetries. A spatial loop
winding once probes a \emph{torelon}, a flux tube wrapped around the box;
a centre transformation distinguishes it from the scalar loop. A pair of
winding tubes can be centre neutral and can enter the scalar channel.
At $\beta=2.5$ and spatial side 10 the single-torelon energy is
$aE_T=0.2368(13)$, so the \emph{noninteracting estimate} for a pair is
$2aE_T\simeq0.474$. A late-time scalar fit gives $0.484(17)$ in a window
chosen after examining the data (Section~\ref{ym:num-glueball}). This
suggests what operator to measure next---a neutral pair alongside the
scalar loop. Agreement is not a measurement of the pair interaction or
of mixing. Likewise the homology of sheets obtained by centre projection
in Section~\ref{ym:num-topology} describes the topology of a gauge-fixed
ensemble; it is not the imposed, gauge-invariant flux pairing of
Eq.~\eqref{ym:guide-flux-pairing}.

Finally, a lattice energy is dimensionless. A bound $(\log2)/n$ gives a
physical bound $(\log2)/(a n)$ only at that spacing. If $a n$ diverges
along a sequence of refinements, the latter lower bound tends to zero.
The missing estimate is a seed at \emph{bounded physical side} $a n$ as
$a\to0$, with matched bare parameters on each pair of boxes. The
small-$\beta$ calculation is at the opposite end of the coupling range.
Even with a uniform seed, one must show that the centred, renormalized
correlations converge to those of the proposed continuum state and that
their vectors test the whole sector orthogonal to the vacuum. A theory can
exist and have positive energy while admitting excitations arbitrarily
close to zero. It is their exclusion, not the existence of the Euclidean
measure or the observation of a flux tube, that Eq.~\eqref{ym:gap-statement}
asks for.
